% !PRINTED 303
\weberpage{303}
%<<CONT
halten sein. Es sind aber nicht alle linearen Congruenzgruppen auch umgekehrt metacyklisch, und diese daraus auszusondern, ist noch ein weiteres Problem, das bis jetzt nur in besonderen Fällen gelöst ist $^{1)}$.

Um aber das Problem wenigstens genauer zu formuliren, fügen wir noch folgende Betrachtungen bei.

Wir bezeichnen die vollständige lineare Congruenzgruppe (6) mit $R$, und sondern daraus die darin enthaltene Gruppe der linearen homogenen Substitutionen

\weq{7}{\begin{aligned}
z_1' &\equiv \alpha_{1,1}z_1 + \alpha_{1,2}z_2 + \cdots + \alpha_{1,k}z_k \\
z_2' &\equiv \alpha_{2,1}z_1 + \alpha_{2,2}z_2 + \cdots + \alpha_{2,k}z_k \quad (\mathrm{mod}\ p) \\
&\cdots\cdots\cdots\cdots\cdots\cdots\cdots\cdots\cdots\cdots \\
z_k' &\equiv \alpha_{k,1}z_1 + \alpha_{k,2}z_2 + \cdots + \alpha_{k,k}z_k
\end{aligned}}

mit nicht verschwindender Determinante ab, die \gesperrt{lineare homogene Congruenzgruppe} genannt, die wir mit $S$ bezeichnen wollen. Ferner ist in (6) die $k$ fach cyklische Gruppe $Q$ enthalten, die aus den Substitutionen

\weq{8}{z_1' = z_1 + \alpha_1, \quad z_2' = z_2 + \alpha_2, \dots, \quad z_k' = z_k + \alpha_k}

besteht.

Bedeutet nun $\sigma$ ein Element aus $S$ und $\gamma$ ein Element aus $Q$, so ergiebt sich aus der Compositionsregel §.\,77, (7) ohne Schwierigkeit, dass sich jede Substitution $\pi$ aus der Gruppe $R$ \gesperrt{auf eine und nur auf eine Weise} in jeder der beiden Formen

\weq{9}{\pi = \sigma\gamma = \gamma'\sigma}

darstellen lässt, worin, wenn $\sigma, \gamma$ durch (7), (8) bestimmt sind, $\gamma'$ die Substitution

\weq{10}{z_1' = z_1 + \alpha_1', \quad z_2' = z_2 + \alpha_2', \dots, \quad z_k' = z_k + \alpha_k'}

bedeutet, wenn die $\alpha_i'$ aus den $\alpha_i$ durch die Substitution $\sigma$ abgeleitet sind.

$Q$ ist ein Normaltheiler von $R$; denn ist $\gamma_1$ ein beliebiges Element aus $Q$, so ist bei mehrmaliger Anwendung von (9):

\[
\sigma\gamma\gamma_1 = \gamma'\gamma_1'\sigma = \gamma'\gamma_1'\sigma\gamma^{-1}\gamma = \gamma'\gamma_1'\gamma'^{-1}\sigma\gamma,
\]

also, wenn

\[
\sigma\gamma = \pi, \quad \gamma'\gamma_1'\gamma'^{-1} = \gamma_1''
\]

gesetzt wird,

\[
\pi\gamma_1 = \gamma_1''\pi,
\]
%CONT>>

\begin{center}
\rule{4cm}{0.4pt}
\end{center}

$^{1)}$ \gesperrt{C. Jordan}, Liouville's Journ. de Math\'em. S\'er. (2), Bd. XIII.
